% GENERATED FILE. Edit canonical lessons, then run npm run generate:books.
% canonical-source-hash: fnv1a64:e0c683568b66836a
% canonical-lessons: GE-C298-verloben, GE-C298-verlobung, GE-C298-braut, GE-C298-braeutigam, GE-C298-brautpaar

\chapter{Die Verlobung}
\label{ch:ge-die-verlobung}
\hlchaptermodality{298}

\begin{chapteropening}
\textbf{By the end of this chapter:} \emph{I can tell how two people decide to marry: to get engaged, an engagement, the bride, the bridegroom and the bride and groom.}

Complete the last lesson of chapter 298: I can tell how two people decide to marry: to get engaged, an engagement, the bride, the bridegroom and the bride and groom.
\end{chapteropening}

\section[verloben]{verloben — to get engaged; mostly reflexive: sich verloben}

\label{lesson:GE-C298-verloben}



\begin{quote}
Before the new one: say the German for the degree, then the German for the move.
\end{quote}

\subsection*{You'll want to know: verloben}
\textbf{verloben} — \textquotedblleft{}to get engaged; mostly reflexive: sich verloben\textquotedblright{}.

\begin{quote}
Sie haben sich im Urlaub verlobt.
\end{quote}

\begin{quote}
\emph{They got engaged on holiday.}
\end{quote}

\emph{Sich mit jemandem verloben}. \emph{Verlobt} — engaged.

\begin{cousinweb}
\emph{verloben} is built on \emph{loben}, which EtymDB traces to Proto-Germanic \emph{lubōną} \textquotedblleft{}to praise, to vow\textquotedblright{}: to make a promise.
\end{cousinweb}

\begin{grammarlens}[title={What you've built}]
One more word for telling, explaining and describing at B1.
\end{grammarlens}

\subsection*{Your turn}
\begin{itemize}
\raggedright
  \item \emph{Say it:} \emph{verloben}
  \item \emph{Say it:} \emph{verloben}, once more
  \item \emph{Say it:} a sentence of your own with \emph{verloben}
  \item \emph{Recall:} say the German for the degree, then the German for the move, then say \emph{verloben} again
\end{itemize}

\begin{tcolorbox}[breakable,colback=teal!4,colframe=teal!35!black,title={Before you move on}]
What does verloben mean? (\textquotedblleft{}To get engaged; mostly reflexive: sich verloben\textquotedblright{}.) Say it once more.
\end{tcolorbox}
\section[die Verlobung]{die Verlobung — the engagement (to marry)}

\label{lesson:GE-C298-verlobung}



\begin{quote}
Before the new one: say the German for the move, then the German for to get engaged.
\end{quote}

\subsection*{You'll want to know: die Verlobung}
\textbf{die Verlobung} — \textquotedblleft{}the engagement (to marry)\textquotedblright{}.

\begin{quote}
Zur Verlobung schenkte er ihr einen Ring.
\end{quote}

\begin{quote}
\emph{For the engagement he gave her a ring.}
\end{quote}

Plural: \emph{die Verlobungen}.

\begin{cousinweb}
\emph{Verlobung} is built on \emph{verloben} \textquotedblleft{}to get engaged\textquotedblright{}, from \emph{loben}, Proto-Germanic \emph{lubōną} \textquotedblleft{}to praise, to vow\textquotedblright{} (EtymDB).
\end{cousinweb}

\begin{grammarlens}[title={What you've built}]
One more word for telling, explaining and describing at B1.
\end{grammarlens}

\subsection*{Your turn}
\begin{itemize}
\raggedright
  \item \emph{Say it:} \emph{die Verlobung}
  \item \emph{Say it:} \emph{die Verlobung}, once more
  \item \emph{Say it:} a sentence of your own with \emph{die Verlobung}
  \item \emph{Recall:} say the German for the move, then the German for to get engaged, then say \emph{die Verlobung} again
\end{itemize}

\begin{tcolorbox}[breakable,colback=teal!4,colframe=teal!35!black,title={Before you move on}]
What does die Verlobung mean? (\textquotedblleft{}The engagement (to marry)\textquotedblright{}.) Say it once more.
\end{tcolorbox}
\section[die Braut]{die Braut — the bride}

\label{lesson:GE-C298-braut}



\begin{quote}
Before the new one: say the German for to get engaged, then the German for the engagement.
\end{quote}

\subsection*{You'll want to know: die Braut}
\textbf{die Braut} — \textquotedblleft{}the bride\textquotedblright{}.

\begin{quote}
Die Braut trug ein langes weißes Kleid.
\end{quote}

\begin{quote}
\emph{The bride wore a long white dress.}
\end{quote}

Plural: \emph{die Bräute}.

\begin{cousinweb}
\emph{Braut} is an inherited Germanic word, and EtymDB lists English \emph{bride} as its cousin.
\end{cousinweb}

\begin{grammarlens}[title={What you've built}]
One more word for telling, explaining and describing at B1.
\end{grammarlens}

\subsection*{Your turn}
\begin{itemize}
\raggedright
  \item \emph{Say it:} \emph{die Braut}
  \item \emph{Say it:} \emph{die Braut}, once more
  \item \emph{Say it:} a sentence of your own with \emph{die Braut}
  \item \emph{Recall:} say the German for to get engaged, then the German for the engagement, then say \emph{die Braut} again
\end{itemize}

\begin{tcolorbox}[breakable,colback=teal!4,colframe=teal!35!black,title={Before you move on}]
What does die Braut mean? (\textquotedblleft{}The bride\textquotedblright{}.) Say it once more.
\end{tcolorbox}
\section[der Bräutigam]{der Bräutigam — the bridegroom}

\label{lesson:GE-C298-braeutigam}



\begin{quote}
Before the new one: say the German for the engagement, then the German for the bride.
\end{quote}

\subsection*{You'll want to know: der Bräutigam}
\textbf{der Bräutigam} — \textquotedblleft{}the bridegroom\textquotedblright{}.

\begin{quote}
Der Bräutigam war sehr nervös.
\end{quote}

\begin{quote}
\emph{The bridegroom was very nervous.}
\end{quote}

Plural: \emph{die Bräutigame}.

\begin{cousinweb}
\emph{Bräutigam} is inherited from Old High German \emph{brūtigomo}, from Proto-Germanic \emph{brūdigumô}, built on the word for \textquotedblleft{}bride\textquotedblright{} and \emph{gumô} \textquotedblleft{}man\textquotedblright{} (EtymDB).
\end{cousinweb}

\begin{grammarlens}[title={What you've built}]
One more word for telling, explaining and describing at B1.
\end{grammarlens}

\subsection*{Your turn}
\begin{itemize}
\raggedright
  \item \emph{Say it:} \emph{der Bräutigam}
  \item \emph{Say it:} \emph{der Bräutigam}, once more
  \item \emph{Say it:} a sentence of your own with \emph{der Bräutigam}
  \item \emph{Recall:} say the German for the engagement, then the German for the bride, then say \emph{der Bräutigam} again
\end{itemize}

\begin{tcolorbox}[breakable,colback=teal!4,colframe=teal!35!black,title={Before you move on}]
What does der Bräutigam mean? (\textquotedblleft{}The bridegroom\textquotedblright{}.) Say it once more.
\end{tcolorbox}
\section[das Brautpaar]{das Brautpaar — the bride and groom, the bridal couple}

\label{lesson:GE-C298-brautpaar}



\begin{quote}
Before the new one: say the German for the bride, then the German for the bridegroom.
\end{quote}

\subsection*{You'll want to know: das Brautpaar}
\textbf{das Brautpaar} — \textquotedblleft{}the bride and groom, the bridal couple\textquotedblright{}.

\begin{quote}
Alle Gäste warteten auf das Brautpaar.
\end{quote}

\begin{quote}
\emph{All the guests waited for the bride and groom.}
\end{quote}

Plural: \emph{die Brautpaare}.

\begin{cousinweb}
\emph{Brautpaar} joins the inherited \emph{Braut} \textquotedblleft{}bride\textquotedblright{} and \emph{Paar} \textquotedblleft{}pair, couple\textquotedblright{}, which came from Latin \emph{pār} \textquotedblleft{}equal\textquotedblright{} (Lewis \& Short); EtymDB lists no parent for the compound.
\end{cousinweb}

\begin{grammarlens}[title={What you've built}]
One more word for telling, explaining and describing at B1.
\end{grammarlens}

\subsection*{Your turn}
\begin{itemize}
\raggedright
  \item \emph{Say it:} \emph{das Brautpaar}
  \item \emph{Say it:} \emph{das Brautpaar}, once more
  \item \emph{Say it:} a sentence of your own with \emph{das Brautpaar}
  \item \emph{Recall:} say the German for the bride, then the German for the bridegroom, then say \emph{das Brautpaar} again
\end{itemize}

\begin{tcolorbox}[breakable,colback=teal!4,colframe=teal!35!black,title={Before you move on}]
What does das Brautpaar mean? (\textquotedblleft{}The bride and groom, the bridal couple\textquotedblright{}.) Say it once more.
\end{tcolorbox}
